%! Author = sriadityavedantam
%! Date = 3/30/26

\section{Properties of Groups}

\begin{definition}[Group under Multiplication]
    Suppose $G$ is a group $(G,*)$, where
    \[
        *:G\times G\to G.
    \]
    Then:
    \begin{gather*}
        1\in G,\\
        1*a = a*1 \quad \forall a\in G,\\
        (a*b)*c = a*(b*c) \quad \forall a,b,c\in G,\\
        \forall a\in G,\ \exists a^{-1}\in G \text{ such that } a^{-1}*a = aa^{-1} = 1.\\
    \end{gather*}
    Also,
    \[
        a^2 = a*a,
    \]
    and inductively,
    \[
        a^n = a(a^{n-1}) = a*a*\dots*a.
    \]
\end{definition}

$S_X$ is the group of permutations on $X$ under composition, and it is noncommutative for $\#X\geq 3$.
If $(\mathcal{R},+,*)$ is a ring, then $(\mathcal{R},+)$ is a group under addition, and it is commutative.
Note that $(\mathcal{R},+,*)$ is ring notation, not group notation.
This means that $0\in \mathcal{R}$ is allowed even though multiplication is closed, because we are not saying $(\mathcal{R},*)$ is a group.

Also,
\[
    U(\mathcal{R}) = \mathcal{R}^* = \{u\in \mathcal{R} : \exists v\in \mathcal{R} \text{ with } uv=vu=1\}
\]
is a group.

\begin{example}
    \[
        \mathbb{Z}_m^*
    \]
    is the group of units in $\mathbb{Z}_m$.
    For example,
    \begin{gather*}
        \mathbb{Z}_6^* = \{1,5\},\\
        \mathbb{Z}_8^* = \{1,3,5,7\},\\
        \mathbb{Z}_p^* = \{1,2,\dots,p-1\}.\\
    \end{gather*}
\end{example}

If $\mathcal{R}$ is any ring, then
\[
    \mat{n}(\mathcal{R})
\]
is the ring of $n\times n$ matrices with entries in $\mathcal{R}$.
Also,
\[
    \mat{n}(\mathcal{R})^*
\]
is the group of invertible matrices with entries in $\mathcal{R}$.
If $n\geq 2$, then $\mat{n}(\mathcal{R})^*$ is nonabelian.
If $\mathbb{F}$ is a field, then matrices represent transformations of vector spaces.
$\mat{n}(\mathbb{F})$ corresponds to linear transformations of an $n$-dimensional vector space over $\mathbb{F}$.

Suppose $A\in \mat{n}(\mathbb{F})$, and $V$ is an $n$-dimensional vector space.
We define
\[
    T(v)=Av,
\]
so $T:V\to V$ is a linear transformation.
Then
\[
    T(\lambda v+\mu w)=\lambda T(v)+\mu T(w).
\]

\begin{definition}[General Linear Group]
    The group of invertible $n\times n$ matrices with entries in $\mathbb{F}$ under multiplication is denoted by
    \[
        GL_n(\mathbb{F}).
    \]
    This is called the General Linear Group.
\end{definition}

\begin{example}
    $GL_2(\mathbb{F}_2)$ consists of matrices
    \[
        \begin{pmatrix}
            a&b\\
            c&d
        \end{pmatrix}
    \]
    with $a,b,c,d\in\{0,1\}$, and such a matrix is invertible if and only if
    \[
        ad-bc\neq 0.
    \]
    The invertible matrices are
    \begin{gather*}
        \begin{pmatrix}
            1&0\\
            0&1
        \end{pmatrix},
        \quad
        \begin{pmatrix}
            1&1\\
            0&1
        \end{pmatrix},
        \quad
        \begin{pmatrix}
            1&1\\
            1&0
        \end{pmatrix},\\
        \begin{pmatrix}
            0&1\\
            1&0
        \end{pmatrix},
        \quad
        \begin{pmatrix}
            0&1\\
            1&1
        \end{pmatrix},
        \quad
        \begin{pmatrix}
            1&0\\
            1&1
        \end{pmatrix}.\\
    \end{gather*}
    So
    \[
        \#GL_2(\mathbb{F}_2)=6,
    \]
    which is isomorphic to $S_3$.
\end{example}

Coming from geometry, we also get groups.
If you take a solid or a regular $n$-gon and consider its rigid motions, then these form a group.
A rigid motion is a motion that picks up the figure and puts it back in place, such as a rotation or reflection.
\[
    \definecolor{zzttqq}{rgb}{0.6,0.2,0}
    \definecolor{ududff}{rgb}{0.30196078431372547,0.30196078431372547,1}
    \begin{tikzpicture}[line cap=round,line join=round,>=triangle 45,x=1cm,y=1cm]
        \clip(-4.810729327044825,-0.9698720753575509) rectangle (3.9701527617352204,3.62171081782759);
        \fill[line width=2pt,color=zzttqq,fill=zzttqq,fill opacity=0.10000000149011612] (-2.66,-0.24) -- (2.3,-0.2) -- (-0.38,3.38) -- cycle;
        \draw [line width=2pt,color=zzttqq] (-2.66,-0.24)-- (2.3,-0.2);
        \draw [line width=2pt,color=zzttqq] (2.3,-0.2)-- (-0.38,3.38);
        \draw [line width=2pt,color=zzttqq] (-0.38,3.38)-- (-2.66,-0.24);
        \begin{scriptsize}
            \draw [fill=ududff] (-2.66,-0.24) circle (2.5pt);
            \draw[color=ududff] (-2.605104916119032,-0.09247745906159882) node {$A$};
            \draw [fill=ududff] (2.3,-0.2) circle (2.5pt);
            \draw[color=ududff] (2.354082045553743,-0.050861904138470665) node {$B$};
            \draw [fill=ududff] (-0.38,3.38) circle (2.5pt);
            \draw[color=ududff] (-0.32318532116750337,3.5280758192505517) node {$C$};
        \end{scriptsize}
    \end{tikzpicture}
\]

\begin{example}
    Rigid motions of a regular $3$-gon are:
    the identity,
    rotation by $120^\circ=\frac{2\pi}{3}$ counterclockwise, call this $r$,
    rotation by $240^\circ=\frac{4\pi}{3}$, call this $r^2$,
    and reflections.
    \begin{gather*}
        1: 1\to 1,\ 2\to 2,\ 3\to 3\\
        r: 1\to 2\to 3\to 1\\
        r^2: 1\to 3\to 2\to 1\\
        s: 1\to 1,\ 2\leftrightarrow 3\\
        sr: 1\leftrightarrow 3,\ 2\to 2\\
        rs: \text{reflection through vertex }3.\\
    \end{gather*}
    This group is denoted $D_3$.
    It has order $6$ and is called the dihedral group.
    Some texts use $D_6$.
    Others use $D_n$ or $D_{2n}$ depending on convention.
\end{example}

\begin{example}
    Let $H,K$ be groups.
    Then $H\times K$ has underlying set $H\times K$ with operation
    \[
        (h_1,k_1)*(h_2,k_2)=(h_1*h_2,k_1*k_2).
    \]
    This makes $H\times K$ into a group.
    The identity is $(1,1)$.
    Associativity holds componentwise.
    The inverse is
    \[
        (h,k)^{-1}=(h^{-1},k^{-1}).
    \]
    Closure also holds componentwise.
\end{example}

$(\mathbb{Z}_2\times\mathbb{Z}_2,+)$ has order $4$.
$(\mathbb{Z}_4,+)$ also has $4$ elements.
$\mathbb{Z}_5^*$ and $\mathbb{Z}_8^*$ also have $4$ elements.
Now, are they isomorphic.
Let us look at the tables for these groups.

\begin{center}
    \begin{tabular}{lllll}
        $(\mathbb{Z}_4,+)$ & 0 & 1 & 2 & 3 \\
        0 & 0 & 1 & 2 & 3 \\
        1 & 1 & 2 & 3 & 0 \\
        2 & 2 & 3 & 0 & 1 \\
        3 & 3 & 0 & 1 & 2
    \end{tabular}

    \vspace{1em}

    \begin{tabular}{lllll}
        $\mathbb{Z}_5^*$ & 1 & 2 & 3 & 4 \\
        1 & 1 & 2 & 3 & 4 \\
        2 & 2 & 4 & 1 & 3 \\
        3 & 3 & 1 & 4 & 2 \\
        4 & 4 & 3 & 2 & 1
    \end{tabular}
\end{center}

A possible isomorphism is
\begin{align*}
    0 &\mapsto 1,\\
    1 &\mapsto 2,\\
    2 &\mapsto 4,\\
    3 &\mapsto 3.
\end{align*}

\begin{theorem}[Fundamental Theorem of Finite Abelian Groups]{
    Every finite abelian group is a direct product of groups of the form $(\mathbb{Z}_n,+)$ for various $n$.
}
\end{theorem}

\begin{definition}[Other Properties of Groups]
    Let $G$ be a group.
    Some other properties which follow from the definition are the following.
    \begin{enumerate}
        \item Cancellation law.
        Suppose $a,b,c\in G$.
        If
        \[
            ba=ca,
        \]
        then
        \[
            b=c.
        \]

        \item Identity is unique.

        \item Inverses are unique.

        \item
        \[
            (ab)^{-1}=b^{-1}a^{-1}.
        \]

        \item
        \[
            (a^{-1})^{-1}=a.
        \]
    \end{enumerate}
\end{definition}

\begin{proposition}{
    The cancellation law holds in any group.
}
    Suppose
    \[
        ab=ac.
    \]
    Then
    \begin{align*}
        a^{-1}(ab) &= a^{-1}(ac),\\
        (a^{-1}a)b &= (a^{-1}a)c,\\
        1b &= 1c,\\
        b &= c.
    \end{align*}
    A similar argument proves right cancellation.
\end{proposition}

\begin{proposition}{
    The identity element in a group is unique.
}
    Suppose $e$ and $f$ are both identities.
    Then
    \[
        ef=f
    \]
    since $e$ is an identity, and
    \[
        ef=e
    \]
    since $f$ is an identity.
    Thus
    \[
        e=f.
    \]
\end{proposition}

\begin{proposition}{
    Inverses in a group are unique.
}
    Suppose $a\in G$, and $b,c\in G$ both satisfy
    \[
        ba=ab=e
        \quad\text{and}\quad
        ca=ac=e.
    \]
    Then
    \[
        b=be=b(ac)=(ba)c=ec=c.
    \]
    Thus $b=c$.
\end{proposition}

\begin{proposition}{
    For all $a,b\in G$,
    \[
        (ab)^{-1}=b^{-1}a^{-1}.
    \]
}
    We compute
    \[
        (ab)(b^{-1}a^{-1})=abb^{-1}a^{-1}=aea^{-1}=aa^{-1}=e,
    \]
    and similarly
    \[
        (b^{-1}a^{-1})(ab)=e.
    \]
    By uniqueness of inverse,
    \[
        (ab)^{-1}=b^{-1}a^{-1}.
    \]
\end{proposition}

\begin{proposition}{
    For all $a\in G$,
    \[
        (a^{-1})^{-1}=a.
    \]
}
    Since
    \[
        a^{-1}a=aa^{-1}=e,
    \]
    $a$ is an inverse of $a^{-1}$.
    By uniqueness of inverse,
    \[
        (a^{-1})^{-1}=a.
    \]
\end{proposition}

\begin{definition}[Order]
    Let $G$ be a group and let $a\in G$.
    The order of $a$ is the smallest strictly positive integer $n$ such that
    \[
        a^n \coloneqq a*a*\dots*a = 1.
    \]
    If there is no such $n$, then $a$ has infinite order.
\end{definition}

\begin{theorem}{
    If $\#G$ is finite, then every $a\in G$ has finite order.
}
    Since $G$ is finite, the list
    \[
        1,a,a^2,a^3,\dots
    \]
    cannot consist of distinct elements forever.
    So there exist integers $i<j$ such that
    \[
        a^i=a^j.
    \]
    By cancellation,
    \[
        1=a^{j-i}.
    \]
    Thus $a$ has finite order.
\end{theorem}

\begin{example}
    Dihedral groups have finite order for all elements.
\end{example}

\begin{example}[$G_1=(\mathbb{Z}_4,+)$]
    \begin{tabular}{ll}
        Elements & Order \\
        0 & 1 \\
        1 & 4 \\
        2 & 2 \\
        3 & 4
    \end{tabular}
\end{example}

\begin{example}[$G_2=\mathbb{Z}_2\times\mathbb{Z}_2$]
    \begin{tabular}{ll}
        Elements & Order \\
        (0,0) & 1 \\
        (0,1) & 2 \\
        (1,0) & 2 \\
        (1,1) & 2
    \end{tabular}
\end{example}

\begin{example}[$G_3=\mathbb{Z}_5^*$]
    \begin{tabular}{ll}
        Elements & Order \\
        1 & 1 \\
        2 & 4 \\
        3 & 4 \\
        4 & 2
    \end{tabular}
\end{example}

This is another way of showing that $G_1$ and $G_3$ are isomorphic.

\begin{corollary}[Infinite Order]{
    Suppose $a^n\neq e$ for all $n>0$.
    Then $|a|=\infty$.
}
\end{corollary}

\begin{theorem}{
    Suppose $a\in G$ and $|a|=n$.
    If $k\in\mathbb{Z}$, then
    \[
        a^k=1
        \quad\text{if and only if}\quad
        n\mid k.
    \]
}
    ($\implies$).
    Suppose $k\in\mathbb{Z}$ such that $a^k=1$.
Write
    \[
        k=nq+r
    \]
    with $0\leq r<n$.
    Then
    \begin{align*}
        a^k &= a^{nq+r}\\
        &= a^{nq}a^r\\
        &= (a^n)^q a^r\\
        &= 1^q a^r\\
        &= a^r.
    \end{align*}
    Since $a^k=1$, we get $a^r=1$.
    But $n$ is the smallest positive integer with this property, so $r=0$.
    Thus $n\mid k$.

    ($\impliedby$).
    If $n\mid k$, then $k=nq$ for some $q$.
    Thus
    \begin{align*}
        a^k &= a^{nq}\\
        &= (a^n)^q\\
        &= 1^q\\
        &= 1.
    \end{align*}
    So $a^k=1$.
\end{theorem}

\begin{corollary}{
    \[
        a^i=a^j
        \quad\text{if and only if}\quad
        a^{i-j}=1.
    \]
    Thus
    \[
        n\mid i-j
    \]
    or equivalently
    \[
        i\equiv j\mod n.
    \]
}
($\implies$).
    Given $a^i=a^j$, then
    \begin{align*}
        a^i a^{-j} &= a^j a^{-j},\\
        a^{i-j} &= 1.
    \end{align*}

    ($\impliedby$).
    Given $a^{i-j}=1$, then
    \begin{align*}
        a^i a^{-j} &= 1,\\
        a^i &= a^j.
    \end{align*}
\end{corollary}

\begin{theorem}{
    Suppose $|a|=n$ and $t\mid n$.
    Then
    \[
        |a^t|=\frac{n}{t}.
    \]
}
    Let $|a^t|=k$.
    Since $a^n=1$, we have
    \[
        (a^t)^{n/t}=a^n=1,
    \]
    so
    \[
        k\mid \frac{n}{t}.
    \]
    Also, since $(a^t)^k=1$, we have
    \[
        a^{tk}=1.
    \]
    Thus
    \[
        n\mid tk.
    \]
    Since $t\mid n$, it follows that
    \[
        \frac{n}{t}\mid k.
    \]
    Therefore
    \[
        k=\frac{n}{t}.
    \]
\end{theorem}

Suppose $|a|=6$, then
\[
    |a^4|=3.
\]

More generally, for $t\in\mathbb{Z}$ with $t>0$, we have
\[
    |a^t|=\frac{n}{\gcd(t,n)}.
\]

\begin{proposition}{
    If $|a|=n$ and $d=\gcd(t,n)$, then
    \[
        |a^t|=\frac{n}{d}.
    \]
}
    Let $|a^t|=k$.
    We know $a^n=1$.
    Then
    \[
        (a^t)^{n/d}=a^{tn/d}=\left(a^n\right)^{t/d}=1,
    \]
    since $t/d\in\mathbb{Z}$.
    So
    \[
        k\mid \frac{n}{d}.
    \]

    Also, since $(a^t)^k=1$, we have
    \[
        a^{tk}=1,
    \]
    so
    \[
        n\mid tk.
    \]
    Dividing by $d$ gives
    \[
        \frac{n}{d}\mid k\cdot \frac{t}{d}.
    \]
    Because $\gcd(n/d,t/d)=1$, it follows that
    \[
        \frac{n}{d}\mid k.
    \]
    Thus
    \[
        k=\frac{n}{d}.
    \]
\end{proposition}